\documentclass[preprint,12pt,authoryear]{elsarticle}
% Neural Networks (Elsevier) -- manuscript source. Build: latexmk -pdf main.tex
\usepackage{amsmath,amssymb,amsthm}
\usepackage{booktabs,graphicx,xcolor}
\usepackage{algorithm}
\usepackage[noend]{algpseudocode}
\usepackage{tikz}
\usetikzlibrary{arrows.meta,positioning,fit,calc}
\usepackage[colorlinks=true,linkcolor=blue!50!black,citecolor=blue!50!black,urlcolor=blue!50!black]{hyperref}
\usepackage[capitalise,nameinlink]{cleveref}
\usepackage{lineno}
\usepackage{float}

\newtheorem{definition}{Definition}
\newtheorem{lemma}{Lemma}
\newtheorem{proposition}{Proposition}
\newtheorem{theorem}{Theorem}
\newtheorem{remark}{Remark}
\crefname{algorithm}{Algorithm}{Algorithms}

% drafting helpers (remove before submission)
\newcommand{\todo}[1]{\textcolor{red!70!black}{\textbf{[TODO:} #1\textbf{]}}}
\newcommand{\pend}{\textcolor{red!70!black}{[pending]}}
\IfFileExists{tables/numbers.tex}{\newcommand{\MeanCertified}{99.4\%}
\newcommand{\MeanViolDomain}{0.02\%}
\newcommand{\FreeViolDomain}{9.7\%}
}{}
\providecommand{\CegisCertified}{\pend}
\providecommand{\CegisUnknown}{\pend}
\providecommand{\MeanCertified}{\pend}
\providecommand{\MeanViolDomain}{\pend}
\providecommand{\FreeViolDomain}{\pend}

\newcommand{\R}{\mathbb{R}}
\newcommand{\B}{\mathcal{B}}
\newcommand{\I}{\mathcal{I}}
\newcommand{\A}{\mathcal{A}}
\newcommand{\sig}{\sigma}
\DeclareMathOperator*{\argmin}{arg\,min}
\DeclareMathOperator*{\argmax}{arg\,max}
\DeclareMathOperator{\clip}{clip}

\journal{Neural Networks}

\begin{document}
\linenumbers

\begin{frontmatter}

\title{Constructive neural networks with certified partial monotonicity for binary classification}

\author[ufba]{Ghabriel A. G. de S\'a\corref{cor}}
\ead{ghabriel\_ags@hotmail.com}
\author[ufba]{Cristiano H. O. Fontes}
\ead{cfontes@ufba.br}
\author[ufba]{Marcelo Embiru\c{c}u}
\ead{embirucu@ufba.br}
\cortext[cor]{Corresponding author.}
\affiliation[ufba]{organization={Programa de Pós-Graduação em Engenharia Industrial, Escola Politécnica, Universidade Federal da Bahia},
  city={Salvador}, state={BA}, country={Brazil}}

\begin{abstract}
Prior knowledge that the probability of an outcome must increase or decrease with some inputs --- risk with age,
malignancy with tumour size, default with debt-to-income ratio --- is common in high-stakes classification, yet
standard networks violate it freely, and most constrained networks either fix the architecture in advance or
restrict every weight to a sign. We study constructive single-hidden-layer networks, which grow their hidden layer
one unit at a time, under \emph{partial} monotonicity: only the inputs with a known direction are constrained. The
gains $\partial p/\partial x_i$ and their gradients with respect to all weights are written in closed form and
imposed as inequality constraints of a sequential quadratic program at every constructive step. We show that the
constraint on the average gain used in earlier work does not imply monotonicity, and we replace it with a
counterexample-guided scheme: point-wise constraints at a growing set of anchors, an adversarial search for
violations over the whole input domain, and a sound branch-and-bound certificate that exploits the structure of the
network --- in input space with interval and mean-value bounds, or in pre-activation space, where the bound is exact
and a linear program discards unreachable regions. Model size is selected by Dynamic Cross-Validation, which advances
one model per inner fold in lock-step. On eight datasets, including the COMPAS, Heart Disease and Loan Defaulter
monotonicity benchmarks, under nested cross-validation against sixteen methods, including the Min-Max, Constrained
Monotonic and Lipschitz Monotonic networks, \CegisCertified{} of the final models are certified monotone on the
entire input domain, against \MeanCertified{} with the average-gain constraint, with \todo{accuracy statement after
the campaign}.
\end{abstract}

\begin{keyword}
Monotonic neural networks \sep Partial monotonicity \sep Constructive neural networks \sep Formal verification \sep
Counterexample-guided learning \sep Dynamic cross-validation
\end{keyword}

\end{frontmatter}

% =====================================================================================================
\section{Introduction}\label{sec:intro}

In many classification problems the direction of the effect of some inputs is known before any data are seen. The
probability of death after heart failure should not decrease as the ejection fraction falls; the probability that
a tumour is malignant should not decrease as its radius grows; the probability of default should not decrease as the
debt-to-income ratio increases. A model that violates such relations may still score well on a test set, but its
predictions cannot be defended to a clinician, a regulator or a customer, and its violations tend to appear exactly
where data are scarce \citep{Sill1997,Cano2019,Liu2020}. \emph{Monotonic} neural networks enforce these relations;
\emph{partially} monotonic networks enforce them only for the inputs whose direction is known and leave the others
free \citep{DanielsVelikova2010,You2017}.

Existing approaches fall into three families. \emph{Soft} approaches penalise violations at sampled points
\citep{SillAbuMostafa1996} and give no guarantee outside the sample. \emph{Structural} approaches make monotonicity a
property of the architecture: non-negative weights \citep{ArcherWang1993,DanielsVelikova2010}, min-max networks
\citep{Sill1997}, lattices \citep{Gupta2016,You2017}, integral networks \citep{Wehenkel2019}, constrained activations
\citep{Runje2023} or Lipschitz-bounded residual networks \citep{Nolte2023}; they are monotone by construction but fix
a specialised topology in advance. \emph{Certified} approaches train an ordinary network and verify it afterwards,
with mixed-integer programming \citep{Liu2020} or SMT solvers \citep{Sivaraman2020}; they are flexible but tied to
piecewise-linear networks and to general-purpose solvers.

\emph{Constructive} networks \citep{Ash1989,Fahlman1990,KwokYeung1997,Lehtokangas1999} address a different question
--- how many hidden units --- by growing the hidden layer until a validation criterion stops improving. Constructive
growth and monotonicity have not been combined, and doing so is not immediate: a structural constraint limits
expressiveness at every size \citep{DanielsVelikova2010}, whereas a constraint imposed only on the training data may
hold at every size on the data and still fail elsewhere. In earlier work we imposed \emph{gain-sign} constraints
--- constraints on the sign of $\partial \hat y/\partial x_i$ --- on constructive networks for regression
\citep{Sa2026}, enforcing the sign of the average finite-difference gain. Here we show that this constraint does not
imply monotonicity and develop a version that does, for binary classification.

\paragraph{Contributions}
\begin{enumerate}
\item \emph{Analytic gain constraints for constructive classifiers.} For a sigmoid single-hidden-layer network the gain
$\partial p/\partial x_i$ and its gradient with respect to every weight have closed forms (\cref{sec:gains}); we impose
them as inequality constraints of a sequential quadratic program \citep{Kraft1988} at each constructive step, with
three constraint families of increasing strength (average, point-wise and sign constraints).
\item \emph{Counterexample-guided training with a sound certificate} (\cref{sec:cegis,sec:certificate}). Point-wise
constraints are imposed at a set of anchors that grows with the violations found by an adversarial search over the
whole input domain; the final model is checked by a branch-and-bound procedure that is sound by construction
(\cref{thm:sound}). Unlike sign constraints, the scheme lets individual hidden units be non-monotone as long as their
combination is, and unlike MILP or SMT verification it needs no external solver: the bounds exploit the fact that the
sign of a gain is the sign of a positive combination of $\sigma'$ terms, which is separable in pre-activation space.
\item \emph{Dynamic Cross-Validation} (DCV, \cref{sec:dcv}): one model per inner fold is grown in lock-step and the size
is chosen by the aggregated validation loss, inside an outer cross-validation loop that is never used for selection.
\item \emph{A rigorous comparison} (\cref{sec:setup}): eight datasets, including the three classification benchmarks of
the monotonic-network literature, nested cross-validation, sixteen methods including Min-Max, Constrained Monotonic and
Lipschitz Monotonic networks, certificate rates in addition to accuracy, and ablations that separate the effect of the
constraints from that of the optimiser. One of them shows that the advantage of weight reuse over random
re-initialisation reported in our earlier constructive work is, to a large extent, an optimiser effect
(\cref{sec:ablation}).
\end{enumerate}

\Cref{sec:related} reviews related work, \cref{sec:problem} fixes notation, \cref{sec:method} presents the method,
\cref{sec:setup} the experimental protocol, \cref{sec:results} the results, \cref{sec:discussion} discusses them and
\cref{sec:conclusion} concludes. Proofs are in \ref{app:proofs}.

% =====================================================================================================
\section{Related work}\label{sec:related}

\paragraph{Monotonic networks}
\citet{ArcherWang1993} constrained back-propagation weights to be non-negative for two-group classification;
\citet{SillAbuMostafa1996} introduced monotonicity hints, penalties on violations at sampled pairs of points; and
\citet{Sill1997} proposed min-max networks, monotone by construction and universal for monotone functions.
\citet{DanielsVelikova2010} proved universal approximation for partially monotone networks with non-negative weights and
sigmoid units in the constrained paths. Lattice models \citep{Gupta2016,You2017} interpolate monotone look-up tables;
UMNN \citep{Wehenkel2019} integrates a positive network. Constrained Monotonic Neural Networks \citep{Runje2023} replace
the single convex activation implied by non-negative weights with a mixture of convex, concave and saturated units,
removing the expressiveness limitation of earlier constructions; Lipschitz Monotonic Networks \citep{Nolte2023} add a
linear monotone term to a network whose Lipschitz constant is bounded by weight normalisation; Scalable Monotonic
Networks \citep{KimLee2024} separate monotone and free paths. For extreme learning machines, \citet{Zhu2017} solved a
quadratic program on the output weights of a fixed random layer. All of these fix the architecture in advance.

\paragraph{Certified monotonicity}
\citet{Liu2020} train ReLU networks with a monotonicity regulariser and certify the result by mixed-integer linear
programming, increasing the penalty until certification succeeds; \citet{Sivaraman2020} (COMET) use an SMT solver to
find counterexamples, augment the training data with them and compute monotone envelopes at prediction time. Both
target piecewise-linear networks and general-purpose solvers. Our certificate is closer in spirit to interval methods
\citep{Moore2009} and interval bound propagation \citep{Gowal2018} used for robustness verification
\citep{Katz2017,Tjeng2019}, specialised to the gain of a smooth single-hidden-layer network, and our training loop to
counterexample-guided inductive synthesis \citep{SolarLezama2008}: a learner that satisfies the constraints on a finite
set and a verifier that either proves them everywhere or returns a point that enters the set.

\paragraph{Constructive networks and model selection}
Constructive algorithms add hidden units one at a time \citep{Ash1989,Fahlman1990,Platt1991,KwokYeung1997}, either
freezing the existing units \citep{Fahlman1990} or retraining all weights \citep{Lehtokangas1999}. Weight
initialisation by linearisation around the optimal linear model, combined with growth and weight reuse (WILCAR), was
proposed for regression and multiclass classification \citep{Sa2022,Fontes2021}; extreme learning machines
\citep{Huang2006} are a common constructive baseline because each size is solved in closed form. Model size has been
chosen with a static hold-out set \citep{Setiono2001} and with cross-validation over candidate architectures
\citep{Aran2009}. Selecting hyper-parameters on the data used to report performance biases the estimate
\citep{VarmaSimon2006,CawleyTalbot2010}; our protocol keeps an outer loop that is never used for selection. Warm
starting from previously trained weights can hurt generalisation \citep{AshAdams2020}, which makes weight reuse in
constructive networks a question to test rather than assume.

% =====================================================================================================
\section{Problem setting}\label{sec:problem}

Let $x\in\R^p$ be scaled by min-max normalisation fitted on the training data, so that the training inputs lie in the
box $\B=[0,1]^p$, and let $y\in\{0,1\}$. A single-hidden-layer feedforward network (SLFN) with $n$ sigmoid units models
\begin{equation}\label{eq:slfn}
p_\theta(x)=\sig\bigl(z_o(x)\bigr),\qquad z_o(x)=\sum_{j=1}^{n} v_j\,\sig\bigl(z_j(x)\bigr)+c,\qquad z_j(x)=w_j^\top x+b_j ,
\end{equation}
with $\theta=(W,b,v,c)$, $W\in\R^{n\times p}$ with rows $w_j^\top$, and $\sig(t)=1/(1+e^{-t})$; $p_\theta(x)$ estimates
$P(Y=1\mid x)$ and the class is assigned at threshold $0.5$.

Prior knowledge is encoded in a \emph{signal vector} $s\in\{-1,0,+1\}^p$: $s_i=+1$ ($-1$) states that $P(Y=1\mid x)$
does not decrease (increase) with $x_i$ when the other inputs are fixed, and $s_i=0$ leaves $x_i$ unconstrained. Let
$\I=\{i: s_i\neq0\}$.

\begin{definition}[Partial monotonicity]\label{def:mono}
$p_\theta$ is \emph{$s$-monotone on $\B$} if, for every $i\in\I$, $x\in\B$ and $t\ge0$ with $x+te_i\in\B$,
$s_i\bigl(p_\theta(x+te_i)-p_\theta(x)\bigr)\ge0$. Since $p_\theta$ is smooth, this holds if and only if the
\emph{gains} $g_i(x;\theta)=\partial p_\theta(x)/\partial x_i$ satisfy $s_i\,g_i(x;\theta)\ge0$ for all $x\in\B$,
$i\in\I$.
\end{definition}

The guarantee is stated on $\B$, the region covered by the training data after scaling; test points outside it are
extrapolations, for which no data-driven model can be certified without further assumptions.

% =====================================================================================================
\section{Method}\label{sec:method}

\subsection{Constructive networks}\label{sec:constructive}

A constructive method trains a sequence of networks $\theta_1,\theta_2,\dots$ with $n=1,2,\dots$ hidden units, stopping
when a validation criterion has not improved for a number of steps (the \emph{patience}) and returning the best size
$n^\ast$. We study three members of this family and their constrained variants.

\emph{WILCAR} \citep{Sa2022,Fontes2021} initialises the first unit by linearising the network around the mean input so
that it matches the least-squares linear model, and \emph{reuses weights}: the network with $n$ units starts from the
trained network with $n-1$ units plus one new unit, so that successive networks form a continuous path.
\emph{RIXM} re-initialises all weights at every size with Glorot initialisation \citep{Glorot2010}; successive networks
are unrelated. \emph{ELM} \citep{Huang2006} draws the hidden layer at random (ReLU units, weights and biases uniform in
$[-1,1]$) and solves the output layer by least squares. The constrained variants are WILCAR-C, RIXM-C and ELM-C.

WILCAR and RIXM minimise the regularised cross-entropy
\begin{equation}\label{eq:loss}
J(\theta)=-\frac1N\sum_{k=1}^{N}\bigl[y_k\log p_\theta(x_k)+(1-y_k)\log\bigl(1-p_\theta(x_k)\bigr)\bigr]
+\frac{\lambda_0}{N}\bigl(\|W\|_F^2+\|v\|_2^2\bigr),
\end{equation}
the negative log-posterior (per sample) under a Gaussian prior of precision $\lambda_0$ on the weights. Unconstrained
networks are trained with L-BFGS \citep{LiuNocedal1989} and constrained ones with SLSQP \citep{Kraft1988}, both
quasi-Newton methods using the analytic gradient of \eqref{eq:loss}, so that differences between a method and its
constrained variant are due to the constraints and not to the optimiser (\cref{sec:ablation}). The logistic output of
ELM-C is fitted by SLSQP on the output weights, with the hidden layer fixed.

\subsection{Closed-form gains}\label{sec:gains}

\begin{lemma}[Gains]\label{lem:gain}
For the network \eqref{eq:slfn},
\begin{equation}\label{eq:gain}
g_i(x;\theta)=\sig'\bigl(z_o(x)\bigr)\sum_{j=1}^{n}v_j\,\sig'\bigl(z_j(x)\bigr)\,W_{ji},
\qquad \sig'(t)=\sig(t)\bigl(1-\sig(t)\bigr),
\end{equation}
and, with $c_{ij}=s_i\,v_j\,W_{ji}$ and $h_i(x;\theta)=\sum_{j=1}^{n}c_{ij}\,\sig'(z_j(x))$,
$s_i\,g_i(x;\theta)=\sig'(z_o(x))\,h_i(x;\theta)$. Since $\sig'>0$, $s_i g_i(x;\theta)\ge0 \iff h_i(x;\theta)\ge0$.
\end{lemma}

The derivatives of $g_i$ with respect to every weight follow by differentiating \eqref{eq:gain} and are given in
\ref{app:gradients}; SLSQP therefore receives exact constraint Jacobians instead of finite differences. The quantity
$h_i$ carries the whole sign information and is used by the certificate.

\subsection{Constraint families}\label{sec:constraints}

At every constructive step the constrained methods solve $\min_\theta J(\theta)$ subject to one of the following sets of
constraints, where $\A\subset\B$ is a finite set of \emph{anchors}:
\begin{align}
\text{(mean)}\quad & \frac{1}{|\A|}\sum_{a\in\A}s_i\,g_i(a;\theta)\ \ge\ \varepsilon_{m}, && i\in\I, \label{eq:mean}\\
\text{(point-wise)}\quad & s_i\,g_i(a;\theta)\ \ge\ \varepsilon_{p}, && i\in\I,\ a\in\A, \label{eq:point}\\
\text{(sign)}\quad & c_{ij}=s_i\,v_j\,W_{ji}\ \ge\ 0, && i\in\I,\ j=1,\dots,n. \label{eq:sign}
\end{align}
The mean constraint \eqref{eq:mean}, with $\A$ the training inputs, is the constraint of our earlier work
\citep{Sa2026}. It does not imply monotonicity: an average of $g_i$ can be positive while $g_i$ is negative on a region
of positive measure, and the finite-sample average says nothing about $\B\setminus\A$. The point-wise constraint
\eqref{eq:point} with $k$-means centres as anchors is stronger but still only covers $\A$. The sign constraint is
sufficient everywhere:

\begin{proposition}[Sign constraints]\label{prop:sign}
If \eqref{eq:sign} holds, $p_\theta$ is $s$-monotone on $\R^p$.
\end{proposition}

\Cref{prop:sign} is the classical non-negative-weight construction \citep{ArcherWang1993,DanielsVelikova2010} written
for partial monotonicity. It forces every hidden unit to be monotone in every constrained input with the prescribed
sign, which excludes, for instance, sums of a bump and a ramp that are monotone overall; it is the price of guaranteeing
monotonicity through the parameters alone. For ELM-C, whose hidden layer is fixed, \eqref{eq:sign} is linear in the
output weights; we draw the hidden weights of constrained inputs with the prescribed sign so that \eqref{eq:sign}
reduces to non-negative output weights (a monotone ELM).

If SLSQP ends at an infeasible point, the step is restarted from a new random initialisation, up to $R$ times; a size
at which no feasible network is found counts as a non-improving step.

\subsection{Counterexample-guided constraints}\label{sec:cegis}

The \emph{cegis} mode keeps the flexibility of \eqref{eq:point} and obtains the guarantee of \eqref{eq:sign} by
growing the anchor set with violations (\cref{alg:cegis}). After each feasible solution, an adversarial search minimises
$h_i(\,\cdot\,;\theta)$ over $\B$ for each $i\in\I$ with a bound-constrained quasi-Newton method \citep{Byrd1995} from
several starting points (random points and the training points with the smallest $h_i$), using the closed-form gradient
$\nabla_x h_i(x)=\sum_j c_{ij}\,\sig''(z_j(x))\,w_j$. Every point with $h_i<0$ becomes an anchor and the network is
re-solved from its current weights. After the final refit, the certificate of \cref{sec:certificate} is run; any
counterexample it returns is added to $\A$ and the model re-solved (``repair'').

\begin{algorithm}[t]
\caption{Constructive step $n$ in the \emph{cegis} mode.}\label{alg:cegis}
\begin{algorithmic}[1]
\Require data $(X,y)$, signal vector $s$, anchors $\A$ ($k$-means centres, plus counterexamples found at earlier sizes),
initial weights $\theta^{(0)}$ (WILCAR: previous network plus one unit; RIXM: random)
\State $\theta\gets$ SLSQP: $\min J$ s.t.\ \eqref{eq:point} on $\A$, from $\theta^{(0)}$; restart from random weights while infeasible (at most $R$ times)
\For{$r=1,\dots,R_{\mathrm{cex}}$}
  \State $C\gets\emptyset$
  \For{$i\in\I$}
    \State $\hat x_i\gets$ multi-start L-BFGS-B on $\min_{x\in\B}h_i(x;\theta)$
    \If{$h_i(\hat x_i;\theta)<0$} $C\gets C\cup\{\hat x_i\}$ \EndIf
  \EndFor
  \If{$C=\emptyset$} \textbf{break} \EndIf
  \State $\A\gets\A\cup C$; \ $\theta\gets$ SLSQP on the enlarged $\A$, warm-started at $\theta$
\EndFor
\State \Return $\theta$, $\A$
\end{algorithmic}
\end{algorithm}

\subsection{A sound certificate}\label{sec:certificate}

By \cref{lem:gain}, $p_\theta$ is $s$-monotone on $\B$ if and only if $\min_{x\in\B}h_i(x;\theta)\ge0$ for every
$i\in\I$. \Cref{alg:bb} decides this by branch and bound, for one $i$ at a time, with $c_j=c_{ij}$ and
$h(x)=\sum_j c_j\sig'(w_j^\top x+b_j)$. If $c_j\ge0$ for every $j$, the root is certified immediately
(\cref{prop:sign}). The bounds rely on two elementary facts: $\sig'$ is even, increasing on $(-\infty,0]$ and
decreasing on $[0,\infty)$; and $|\sig''|$ is even, increasing on $[0,t^\ast]$ and decreasing on $[t^\ast,\infty)$, with
$t^\ast=\ln(2+\sqrt3)$ and $|\sig''(t^\ast)|=\sqrt3/18$.

\paragraph{Input-space bounds}
For a box $Q=[l,u]\subseteq\B$ with centre $m$ and radius $r=(u-l)/2$, each pre-activation ranges exactly over
$[\alpha_j,\beta_j]=[w_j^\top m+b_j-|w_j|^\top r,\ w_j^\top m+b_j+|w_j|^\top r]$. With
$\underline{d}_j=\min\{\sig'(\alpha_j),\sig'(\beta_j)\}$, $\overline{d}_j=\sig'(\clip(0;\alpha_j,\beta_j))$ and
$\kappa_j=\max_{t\in[\alpha_j,\beta_j]}|\sig''(t)|$ (equal to $\sqrt3/18$ if $\pm t^\ast\in[\alpha_j,\beta_j]$ and to
the larger endpoint value otherwise),
\begin{equation}\label{eq:lbx}
L(Q)=\max\Bigl\{\sum_{c_j\ge0}c_j\underline{d}_j+\sum_{c_j<0}c_j\overline{d}_j,\ \ h(m)-\sum_{d=1}^{p}r_d\sum_{j}|c_j|\,\kappa_j\,|W_{jd}|\Bigr\}.
\end{equation}
The first term is the natural interval extension \citep{Moore2009}; the second is a mean-value form, which becomes
tight as boxes shrink.

\paragraph{Pre-activation-space bounds}
When $n<p$, $h=\varphi\circ(Wx+b)$ with $\varphi(z)=\sum_j c_j\sig'(z_j)$ \emph{separable}. On a box
$R=[\zeta^l,\zeta^u]\subset\R^n$ its minimum is therefore exact:
\begin{equation}\label{eq:lbz}
\min_{z\in R}\varphi(z)=\sum_{c_j\ge0}c_j\min\{\sig'(\zeta^l_j),\sig'(\zeta^u_j)\}+\sum_{c_j<0}c_j\,\sig'\bigl(\clip(0;\zeta^l_j,\zeta^u_j)\bigr).
\end{equation}
The pre-activations reachable from $\B$ form the zonotope $Z=\{Wx+b:x\in\B\}$; a box with $R\cap Z=\emptyset$ can be
discarded, and this is decided by a linear program in $x$. The search starts from the bounding box of $Z$, splits the
box with the most negative bound first, and, for boxes that meet $Z$, evaluates $h$ at the solution of the linear
program (minimising the linearisation of $\varphi$ at the centre of $R$), which yields counterexamples.

\begin{algorithm}[t]
\caption{Branch-and-bound certificate for input $i$ ($h=h_i$).}\label{alg:bb}
\begin{algorithmic}[1]
\If{$c_j\ge0$ for all $j$} \Return \textsc{certified} \EndIf
\State $\mathcal{Q}\gets\{\B\}$ \Comment{input space; pre-activation space when $n<p$ and the input-space search exhausts $N/4$ nodes}
\While{$\mathcal{Q}\neq\emptyset$}
  \State remove a box $Q$ from $\mathcal{Q}$
  \If{$L(Q)\ge0$ (\eqref{eq:lbx}; \eqref{eq:lbz} in pre-activation space) \textbf{or} $Q$ does not meet $Z$} \textbf{continue} \EndIf
  \State $\hat x\gets$ centre of $Q$ (pre-activation space: point returned by the linear program)
  \If{$h(\hat x)<0$} \Return \textsc{counterexample} $\hat x$ \EndIf
  \If{more than $N$ boxes processed} \Return \textsc{unknown} \EndIf
  \State split $Q$ in half along $\argmax_d (u_d-l_d)\sum_j|c_j||W_{jd}|$ (pre-activation space: $\argmax_j(\zeta^u_j-\zeta^l_j)|c_j|$); add both halves to $\mathcal{Q}$
\EndWhile
\State \Return \textsc{certified}
\end{algorithmic}
\end{algorithm}

\begin{theorem}[Soundness]\label{thm:sound}
If \cref{alg:bb} returns \textsc{certified} for every $i\in\I$, then $p_\theta$ is $s$-monotone on $\B$. If it returns
\textsc{counterexample} $\hat x$, then $\hat x\in\B$ and $s_i\,g_i(\hat x;\theta)<0$.
\end{theorem}

\begin{proposition}[Termination]\label{prop:term}
If $\min_{\B}h_i\neq0$ and boxes are processed in breadth-first order, the input-space search terminates after
finitely many boxes when $N=\infty$.
\end{proposition}

The only undecided case is a minimum exactly at zero, where the network is monotone but touches the boundary; the
search then stops at the budget with \textsc{unknown}, which we report separately. The certificate is exact up to
floating-point rounding; directed rounding \citep{Moore2009} would make it rigorous in that sense too and is left
for an implementation concerned with adversarial use. \Cref{thm:sound} holds for any SLFN, so we also certify the
networks trained with the mean and point-wise constraints and the unconstrained ones, which measures how often their
monotonicity holds without a guarantee.

\subsection{Dynamic cross-validation}\label{sec:dcv}

\Cref{alg:dcv} selects the number of units inside a training set. The inner training set is split into $K$ stratified
folds, one model per fold is created, and all $K$ models are grown in lock-step; at size $n$ the criterion is the sum of
the $K$ validation cross-entropies. Growth stops when the criterion has not improved for the patience, and the final
network is retrained on the whole training set with the selected size $n^\ast$ (for WILCAR, along the constructive
path $1,\dots,n^\ast$). The criterion is informative only if the $K$ curves are smooth functions of $n$, which weight
reuse provides and random re-initialisation does not; we test this by also running DCV for RIXM
(\cref{sec:ablation}). The alternative selection is a stratified 80/20 hold-out split with the validation MCC as
criterion.

\begin{algorithm}[t]
\caption{Dynamic cross-validation (DCV) inside a training set $T$.}\label{alg:dcv}
\begin{algorithmic}[1]
\State split $T$ into $K$ stratified folds $(T_k^{\mathrm{tr}},T_k^{\mathrm{va}})$; initialise one model $M_k$ per fold on $T_k^{\mathrm{tr}}$
\State $J^\ast\gets\infty$, $w\gets0$
\For{$n=1,2,\dots,n_{\max}$}
  \For{$k=1,\dots,K$} \State $M_k\gets$ constructive step $n$ of $M_k$ on $T_k^{\mathrm{tr}}$ (warm start from $M_k$ at $n-1$ if the method reuses weights) \EndFor
  \State $J(n)\gets\sum_k \mathrm{BCE}\bigl(M_k;T_k^{\mathrm{va}}\bigr)$ \ ($J(n)=\infty$ if some $M_k$ is infeasible)
  \If{$J(n)<J^\ast$} $J^\ast\gets J(n)$, $n^\ast\gets n$, $w\gets0$ \Else\ $w\gets w+1$; \textbf{if} $w=$ patience \textbf{break} \EndIf
\EndFor
\State \Return network retrained on $T$ with $n^\ast$ units
\end{algorithmic}
\end{algorithm}

% =====================================================================================================
\section{Experimental setup}\label{sec:setup}

\subsection{Datasets and signal vectors}
\Cref{tab:datasets} lists the eight datasets. The first five come from the UCI repository \citep{UCI}; their signal
vectors, derived from domain knowledge, are given in \ref{app:signals}. The last three are the classification
benchmarks of the monotonic-network literature \citep{Liu2020,Sivaraman2020,Nolte2023,Runje2023}, used with the
preprocessing, train/test split and monotone features of their public release \citep{Runje2023data}. COMPAS and Heart
Disease are small enough for full nested cross-validation; Loan Defaulter is subsampled (stratified) to 20\,000 rows for
cross-validation, because each constrained step solves a dense quadratic program.

\IfFileExists{tables/datasets.tex}{\begin{table}[!tbp]\centering\small
\caption{Datasets. $n$: samples used for cross-validation; $p$: inputs; constrained: inputs with a prescribed gain sign $s_i\neq0$ (constraint ratio); positive: share of the positive class. The last three are the monotonicity benchmarks of \citet{Liu2020,Sivaraman2020,Runje2023}, with their preprocessing and monotone features.}
\label{tab:datasets}
\resizebox{\textwidth}{!}{\begin{tabular}{llllrrr}\toprule
ID & Dataset & Domain & $n$ & $p$ & Constrained & Positive (\%) \\\midrule
ALG & Algerian forest fires \citet{Abid2020} & Fire ecology & 243 & 10 & 9 (90\%) & 56.4 \\
BCO & Breast cancer (original) \citet{Wolberg1990} & Oncology & 683 & 9 & 9 (100\%) & 35.0 \\
BCD & Breast cancer (diagnostic) \citet{Street1993} & Oncology & 569 & 30 & 27 (90\%) & 37.3 \\
HF & Heart failure records \citet{Ahmad2017} & Cardiology & 299 & 12 & 7 (58\%) & 32.1 \\
PIM & Pima Indians diabetes \citet{Smith1988} & Endocrinology & 768 & 8 & 6 (75\%) & 34.9 \\
HD & Heart disease (Cleveland) \citet{Detrano1989} & Cardiology & 303 & 13 & 2 (15\%) & 27.4 \\
COM & COMPAS recidivism \citet{Dressel2018} & Criminal justice & 6172 & 13 & 4 (31\%) & 45.5 \\
LOAN & Loan defaulter \citet{Liu2020} & Credit risk & 20000\textsuperscript{a} & 28 & 5 (18\%) & \pend \\
\bottomrule\end{tabular}}
\par\smallskip\footnotesize \textsuperscript{a}Stratified subsample of 20\,000 rows, used for cross-validation because of the cost of the constructive SLSQP fits (Section~\ref{sec:lit} keeps the official test set whole).
\end{table}}{}

\subsection{Protocol}
Performance is estimated by nested cross-validation \citep{VarmaSimon2006,CawleyTalbot2010}: an outer stratified 5-fold
split repeated five times with different seeds (twice for COMPAS and Loan Defaulter), the outer test fold used exactly
once. Inside each outer training set, inputs are scaled to $[0,1]$ with statistics of that set, the model size (or the
hyper-parameters of the baselines) is selected on a stratified 80/20 split or by DCV with $K=4$, and the final model is
refitted on the whole outer training set. All methods see the same outer and inner splits. For the three benchmarks we
additionally train on the official training set and test on the official test set with five seeds (\cref{sec:lit}),
which is the protocol of the published results.

\subsection{Methods and hyper-parameters}
Constructive networks: WILCAR, RIXM and ELM (unconstrained; $n_{\max}=300$, patience 30) and WILCAR-C, RIXM-C and ELM-C
($n_{\max}=50$, patience 15) with the mean, point-wise, sign and cegis constraints; $\varepsilon_m=0.05$,
$\varepsilon_p=10^{-3}$, ten $k$-means anchors, $R=100$ restarts, $\lambda_0=0.3$, at most 1000 iterations of SLSQP
and L-BFGS, $R_{\mathrm{cex}}=10$ counterexample rounds with 16 starts, and certificate budgets of $2\times10^5$ boxes
(input space) and 3000 linear programs (pre-activation space). Each selection loop has a budget of one hour.
Baselines, with hyper-parameters chosen on the same validation split and refitted on training plus validation data:
logistic regression, free and with sign-bounded coefficients (LR, LR-C); XGBoost \citep{ChenGuestrin2016} and LightGBM
\citep{Ke2017}, free and with monotone constraints (XGB, XGB-C, LGBM, LGBM-C); a one-hidden-layer MLP
\citep{Pedregosa2011}; Min-Max networks \citep{Sill1997}; Constrained Monotonic Neural Networks \citep{Runje2023} and
Lipschitz Monotonic Networks \citep{Nolte2023}, both from their authors' packages. The last three are monotone by
construction and trained with Adam \citep{KingmaBa2015} and early stopping. The grids are listed in
\ref{app:hyper}.

\subsection{Metrics and statistics}
The primary metric is the Matthews correlation coefficient \citep{Matthews1975,ChiccoJurman2020}; we also report the
area under the ROC curve and the Brier score \citep{Brier1950}. Monotonicity is measured by (i) the share of
(point, constrained input) pairs whose finite-difference gain has the wrong sign, on the test points and on $10^4$
uniform points of $\B$, and (ii) the result of \cref{alg:bb} on the final model (certified, counterexample or unknown).
Methods are compared over datasets with the Friedman test, the Iman--Davenport correction and the Nemenyi critical
difference \citep{Demsar2006,ImanDavenport1980}, with Holm's correction \citep{Holm1979} for comparisons with the best
method, and paired within datasets with the Wilcoxon signed-rank test \citep{Wilcoxon1945} and the corrected resampled
$t$-test of \citet{NadeauBengio2003}, which accounts for the overlap between
repeated cross-validation folds.

\subsection{Implementation}
Everything is implemented in Python with NumPy, SciPy \citep{Virtanen2020} (SLSQP, L-BFGS-B, HiGHS linear programs) and
scikit-learn \citep{Pedregosa2011}; the baselines use PyTorch and the packages \texttt{mononet} and
\texttt{monotonicnetworks}. Experiments ran on a 16-thread workstation, one task per core. Code, signal vectors and
raw results are available at \todo{public repository URL}.

% =====================================================================================================
\section{Results}\label{sec:results}

\subsection{Predictive performance}\label{sec:perf}
\IfFileExists{tables/main_mcc.tex}{\begin{table*}[!tbp]\centering
\caption{Test MCC on the real datasets, mean over the outer folds ($5\times5$-fold nested cross-validation; $5\times2$ for LOAN). Constraints as in \cref{sec:problem}; monotone: monotone splits of the tree ensembles; structural: monotone by architecture. Selection of the number of hidden units by DCV or on a hold-out split. Best value per dataset in bold; rank: average rank over the eight datasets. All methods and standard deviations: Supplementary \cref{S-tab:main_mcc_full}.}
\label{tab:main_mcc}
\resizebox{\linewidth}{!}{\begin{tabular}{lllccccccccc}\toprule
Model & Constraint & Selection & ALG & BCO & BCD & HF & PIM & HD & COM & LOAN & Rank \\\midrule
\multicolumn{12}{l}{\textit{Certified SFNNs (this work)}} \\
WILCAR & cegis & DCV & 0.901 & 0.932 & 0.952 & 0.574 & 0.473 & 0.589 & 0.349 & 0.293 & 5.8 \\
RIXM & cegis & hold-out & 0.908 & \textbf{0.936} & 0.937 & 0.564 & 0.471 & 0.599 & 0.352 & 0.292 & 5.2 \\
\midrule
\multicolumn{12}{l}{\textit{Other SFNNs}} \\
WILCAR & sign & DCV & 0.899 & 0.929 & 0.953 & 0.574 & 0.471 & 0.589 & 0.350 & 0.292 & 6.4 \\
WILCAR & mean & DCV & 0.901 & 0.936 & 0.938 & 0.558 & \textbf{0.475} & 0.589 & 0.350 & 0.293 & 5.5 \\
WILCAR & none & DCV & 0.897 & 0.929 & 0.952 & 0.567 & 0.472 & 0.589 & 0.349 & 0.292 & 6.8 \\
RIXM & none & hold-out & 0.907 & 0.931 & 0.954 & 0.563 & 0.473 & 0.599 & 0.351 & 0.294 & 4.5 \\
MLP & none & hold-out & 0.868 & 0.934 & 0.949 & 0.556 & 0.427 & 0.585 & 0.344 & 0.291 & 9.5 \\
\midrule
\multicolumn{12}{l}{\textit{Monotone baselines}} \\
LR & sign & hold-out & 0.931 & 0.924 & \textbf{0.954} & 0.576 & 0.455 & 0.569 & 0.344 & 0.286 & 8.1 \\
XGBoost & monotone & hold-out & \textbf{0.967} & 0.926 & 0.931 & 0.617 & 0.462 & 0.515 & \textbf{0.357} & \textbf{0.295} & 5.9 \\
LightGBM & monotone & hold-out & 0.950 & 0.922 & 0.933 & \textbf{0.628} & 0.459 & 0.518 & 0.356 & 0.295 & 6.4 \\
Min-Max & structural & hold-out & 0.911 & 0.925 & 0.914 & 0.539 & 0.455 & 0.544 & 0.348 & 0.288 & 10.8 \\
CMNN & structural & hold-out & 0.937 & 0.927 & 0.947 & 0.562 & 0.466 & \textbf{0.610} & 0.354 & 0.293 & 5.6 \\
LMN & structural & hold-out & 0.894 & 0.932 & 0.929 & 0.545 & 0.448 & 0.548 & 0.346 & 0.291 & 10.5 \\
\bottomrule\end{tabular}}
\end{table*}}{}
\todo{Describe \cref{tab:main_mcc}: Friedman/Iman--Davenport result, CD diagram (\cref{fig:cd}), where the certified
networks rank relative to the free SLFNs and to CMNN/LMN/Min-Max/XGB-C/LGBM-C; AUC and Brier tables in the
supplementary material.}

\begin{figure}[t]\centering
\IfFileExists{../results/revision/main/cd_diagram.png}{\includegraphics[width=0.8\textwidth]{../results/revision/main/cd_diagram.png}}{\fbox{\parbox{0.8\textwidth}{\centering\vspace{2em}critical-difference diagram \pend\vspace{2em}}}}
\caption{Average MCC ranks over the eight datasets and Nemenyi critical difference ($\alpha=0.05$).}\label{fig:cd}
\end{figure}

\subsection{Monotonicity and certification}\label{sec:cert}
\IfFileExists{tables/certification.tex}{\begin{table*}[!tbp]\centering\small
\caption{Monotonicity on the real data by constraint, pooled over datasets and outer folds. $D$: datasets with results so far; MCC: mean over those datasets; viol.: share of (point, constrained input) pairs whose gain has the wrong sign, on the test points and on $10^4$ uniform points of $[0,1]^p$; cert.: share of final models proved monotone on $[0,1]^p$ by \cref{alg:bb} or monotone by construction (last row: LR-C, XGBoost and LightGBM with monotone constraints, Min-Max, CMNN and LMN, pooled); fallb.: cegis networks returned by the sign-constrained fallback of \cref{prop:always}; $n^\ast$: selected hidden units; ``--'': not applicable or not run yet. Networks selected by DCV behave as those selected on the hold-out split (Supplementary \cref{S-tab:size}).}
\label{tab:cert}
\resizebox{\linewidth}{!}{\begin{tabular}{lllrrrrrrr}\toprule
Constraint & Model & Selection & $D$ & MCC & viol.\ test (\%) & viol.\ domain (\%) & cert.\ (\%) & fallb.\ (\%) & $n^\ast$ \\\midrule
none & WILCAR & hold-out & 3 & 0.930 & 10.96 & 10.96 & 33 & -- & 1.4 \\
 & RIXM & hold-out & 2 & 0.919 & 7.56 & 7.56 & 50 & -- & 1.6 \\
 & MLP & hold-out & 2 & 0.901 & 6.86 & 6.78 & -- & -- & -- \\
\midrule
mean & WILCAR & hold-out & 3 & 0.911 & 0.02 & 0.04 & 99 & -- & 1.5 \\
 & RIXM & hold-out & 2 & 0.920 & 0.00 & 0.00 & 100 & -- & 1.4 \\
\midrule
anchor & WILCAR & hold-out & -- & -- & -- & -- & -- & -- & -- \\
 & RIXM & hold-out & -- & -- & -- & -- & -- & -- & -- \\
\midrule
sign & WILCAR & hold-out & -- & -- & -- & -- & -- & -- & -- \\
 & RIXM & hold-out & -- & -- & -- & -- & -- & -- & -- \\
\midrule
cegis & WILCAR & hold-out & -- & -- & -- & -- & -- & -- & -- \\
 & WILCAR & DCV-1SE & -- & -- & -- & -- & -- & -- & -- \\
 & RIXM & hold-out & -- & -- & -- & -- & -- & -- & -- \\
\midrule
by construction & six monotone baselines & hold-out & 2 & 0.929 & 0.00 & 0.00 & 100 & -- & -- \\
\bottomrule\end{tabular}}\end{table*}}{}
\todo{Describe \cref{tab:cert}: violation rates of free networks (\FreeViolDomain{} of domain points), residual
violations and certification rate of the mean mode (\MeanCertified{}), point-wise mode, sign mode (certified by
construction, accuracy cost) and cegis (\CegisCertified{} certified, \CegisUnknown{} unknown); number of counterexample
rounds and anchors; certificate time.}

\subsection{Comparison with published results}\label{sec:lit}
\IfFileExists{tables/literature.tex}{\begin{table*}[!tbp]\centering\small
\caption{Test accuracy (\%) on the official train/test splits of the benchmarks, mean (standard deviation). Upper block: values reported in the cited papers, each under its own protocol (e.g.\ best five of ten runs for CMNN). Lower block: this work, five seeds, one common protocol (hyper-parameters selected on a validation split of the official training set); ``our runs'': the published architectures trained under that protocol.}
\label{tab:lit}
\resizebox{\linewidth}{!}{\begin{tabular}{lccc}\toprule
Method & COMPAS & Heart Disease & Loan Defaulter \\\midrule
XGBoost \citep{Runje2023} & 68.5 (0.1) & -- & 63.7 (0.1) \\
Min-Max Net \citep{Runje2023} & 67.8 (0.1) & 75 (4) & 64.9 (0.1) \\
DLN \citep{Runje2023} & 67.9 (0.3) & 86 (2) & 65.1 (0.2) \\
Certified MNN \citep{Liu2020} & 68.8 (0.2) & -- & 65.2 (0.1) \\
COMET \citep{Sivaraman2020} & -- & 86 (3) & -- \\
LMN \citep{Nolte2023} & 69.3 (0.1) & 89.6 (1.9) & 65.44 (0.03) \\
CMNN \citep{Runje2023} & 69.2 (0.2) & 89 (0) & 65.3 (0.01) \\
Activation-switch MNN \citep{Sartor2025} & 69.5 (0.1) & 94 (1) & 65.4 (0.1) \\
\midrule
WILCAR, cegis, DCV & 67.9 (0.2) & 82.0 (0.0) & 65.0 (0.0) \\
RIXM, cegis, hold-out & 68.2 (0.3) & 80.3 (0.0) & 65.0 (0.0) \\
WILCAR, unconstrained, DCV & 67.8 (0.1) & 82.0 (0.0) & 65.1 (0.0) \\
LR, sign & 68.1 (0.0) & 82.0 (0.0) & 64.7 (0.1) \\
XGBoost, monotone & 68.6 (0.4) & 85.6 (0.7) & 65.2 (0.1) \\
LightGBM, monotone & 68.8 (0.3) & 85.9 (2.2) & 65.1 (0.1) \\
Min-Max (our runs) & 67.6 (0.5) & 80.3 (2.3) & 64.8 (0.1) \\
CMNN (our runs) & 67.8 (0.2) & 80.7 (0.7) & 65.0 (0.1) \\
LMN (our runs) & 67.7 (0.6) & 83.6 (1.6) & 64.7 (0.5) \\
\bottomrule\end{tabular}}

\end{table*}}{}
\todo{Describe \cref{tab:lit}. Note that published numbers are best-five-of-ten runs with tuned hyper-parameters;
ours are all five seeds.}

\subsection{Ablations}\label{sec:ablation}
\IfFileExists{tables/ablation.tex}{\begin{table}[H]\centering\small
\caption{Ablations of the training recipe on the five UCI datasets (25 outer folds each; constrained methods in the mean mode). MCC: mean over datasets; time: median per outer fold, in seconds. Reference: the same method and selection in the main campaign (L-BFGS, $\lambda_0=0.3$, new unit with random output weight); for the last block, the hold-out selection of the same method.}
\label{tab:ablation}
\begin{tabular}{lllrrrr}\toprule
 & & & \multicolumn{2}{c}{Reference} & \multicolumn{2}{c}{Ablation} \\\cmidrule(lr){4-5}\cmidrule(lr){6-7}
Model & Constraint & Selection & MCC & time & MCC & time \\\midrule
\multicolumn{7}{l}{\emph{No weight prior ($\lambda_0=0$)}}\\
WILCAR & mean & DCV & 0.762 & 99 & 0.661 & 436 \\
WILCAR & mean & hold-out & 0.753 & 7 & 0.633 & 113 \\
RIXM & mean & hold-out & 0.760 & 6 & 0.608 & 337 \\
WILCAR & none & DCV & 0.764 & 17 & 0.755 & 4 \\
WILCAR & none & hold-out & 0.764 & 1 & 0.737 & 1 \\
RIXM & none & hold-out & 0.766 & 4 & 0.726 & 4 \\
\multicolumn{7}{l}{\emph{Gradient descent, $\lambda_0=0$ \citep{Sa2022}}}\\
WILCAR & none & DCV & 0.764 & 17 & 0.775 & 181 \\
WILCAR & none & hold-out & 0.764 & 1 & 0.765 & 40 \\
RIXM & none & hold-out & 0.766 & 4 & 0.672 & 64 \\
\multicolumn{7}{l}{\emph{New unit with zero output weight}}\\
WILCAR & mean & DCV & 0.762 & 99 & 0.762 & 104 \\
WILCAR & mean & hold-out & 0.753 & 7 & 0.753 & 8 \\
WILCAR & none & DCV & 0.764 & 17 & 0.764 & 10 \\
WILCAR & none & hold-out & 0.764 & 1 & 0.765 & 1 \\
\multicolumn{7}{l}{\emph{DCV with re-initialisation (reference: hold-out)}}\\
RIXM & mean & DCV & 0.760 & 6 & 0.761 & 77 \\
RIXM & none & DCV & 0.766 & 4 & 0.764 & 144 \\
\bottomrule\end{tabular}\end{table}}{}
\todo{Describe \cref{tab:ablation}: (i) optimiser: with full-batch gradient descent (dissertation setting) RIXM is
under-trained and WILCAR appears superior; with L-BFGS the gap closes -- report numbers; (ii) prior precision
$\lambda_0$; (iii) zero initialisation of the new unit; (iv) DCV without weight reuse.}

\subsection{Model size and cost}\label{sec:cost}
\todo{$n^\ast$ by method and selection; time per outer fold; WILCAR vs RIXM cost (warm start); cost of cegis
rounds and of the certificate.}

% =====================================================================================================
\section{Discussion}\label{sec:discussion}

\paragraph{What the mean-gain constraint guarantees}
Constraining the average gain over the training data is attractive because it is a single inequality per input, but
\cref{tab:cert} shows that it leaves point-wise violations that the certificate detects. The failure is structural:
an average over a finite sample constrains neither the sign of the gain at each sample point nor its sign between
them. Any claim of ``physical consistency'' must rest on a guarantee over the domain, which only the sign and cegis
modes provide. \todo{quantify with the final numbers.}

\paragraph{Sign constraints versus certified flexibility}
\todo{Compare sign (global) and cegis: accuracy, $n^\ast$, cost; interpret in terms of the hidden units that cegis
lets be non-monotone.}

\paragraph{Weight reuse}
\todo{Interpret the optimiser ablation: warm starting accumulates optimisation across sizes and compensates for a weak
optimiser; with a quasi-Newton optimiser its benefit is cost, smooth selection curves (DCV) and stable constrained
growth, not accuracy. Relate to \citet{AshAdams2020}.}

\paragraph{Limitations}
The certificate is specific to single-hidden-layer networks: separability of $\varphi$ in pre-activation space does not
survive a second hidden layer, where the same scheme would need bound propagation through the layers
\citep{Gowal2018}. Branch and bound is exponential in the worst case; the pre-activation formulation reduces the
dimension from $p$ to $n$, which is small for the selected networks, and undecided cases are reported rather than
hidden. The guarantee is on the training box $\B$, categorical inputs encoded as indicators are treated as
continuous (a conservative choice), and the quadratic programs limit the constrained networks to tens of units and
the training sets to tens of thousands of rows.

% =====================================================================================================
\section{Conclusions}\label{sec:conclusion}
We combined constructive growth of single-hidden-layer networks with partial monotonicity constraints written in
closed form, replaced the average-gain constraint of earlier work --- which does not imply monotonicity --- with a
counterexample-guided scheme whose output is certified on the whole input domain by a solver-free branch-and-bound
procedure, and selected model size with dynamic cross-validation inside a nested protocol. \todo{Summarise the
results.} Extensions to deeper constructive networks and to multiclass problems with monotone class orderings are
natural next steps.

% =====================================================================================================
\section*{CRediT authorship contribution statement}
\todo{confirm.} \textbf{Ghabriel A. G. de S\'a:} Conceptualization, Methodology, Software, Investigation, Formal
analysis, Writing -- original draft. \textbf{Cristiano H. O. Fontes:} Conceptualization, Supervision, Writing --
review \& editing. \textbf{Marcelo Embiru\c{c}u:} Supervision, Writing -- review \& editing.

\section*{Declaration of competing interest}
The authors declare that they have no known competing financial interests or personal relationships that could have
appeared to influence the work reported in this paper.

\section*{Declaration of generative AI and AI-assisted technologies in the writing process}
\todo{confirm wording.} During the preparation of this work the authors used Claude (Anthropic) to assist with code
development and with drafting and editing the text. After using this tool, the authors reviewed and edited the content
as needed and take full responsibility for the content of the publication.

\section*{Data availability}
All datasets are public (\cref{tab:datasets}). Code and results: \todo{repository URL}.

\section*{Funding}
\todo{funding statement or ``This research did not receive any specific grant from funding agencies in the public,
commercial, or not-for-profit sectors.''}

\appendix
% =====================================================================================================
\section{Proofs}\label{app:proofs}

\paragraph{\Cref{lem:gain}}
By the chain rule,
\[
\frac{\partial p_\theta}{\partial x_i}=\sig'(z_o)\,\frac{\partial z_o}{\partial x_i},\qquad
\frac{\partial z_o}{\partial x_i}=\sum_j v_j\sig'(z_j)W_{ji}.
\]
Multiplying by $s_i$ gives $s_ig_i=\sig'(z_o)\sum_j s_iv_jW_{ji}\sig'(z_j)=\sig'(z_o)h_i$, and $\sig'(z_o)>0$.

\paragraph{\Cref{prop:sign}}
If $c_{ij}\ge0$ for all $j$, $h_i(x)=\sum_jc_{ij}\sig'(z_j(x))\ge0$ for every $x\in\R^p$ because $\sig'>0$; conclude
with \cref{lem:gain} and \cref{def:mono}.

\paragraph{Bounds}
(a) For $x\in Q$, $w_j^\top x+b_j=w_j^\top m+b_j+w_j^\top(x-m)$ with $|w_j^\top(x-m)|\le|w_j|^\top r$, attained at a
vertex, so $z_j(Q)=[\alpha_j,\beta_j]$. Since $\sig'$ is even and unimodal with its maximum at $0$, its minimum over
an interval is attained at an endpoint and its maximum at the point of the interval closest to $0$. Hence
$c_j\sig'(z_j(x))\ge c_j\underline d_j$ when $c_j\ge0$ and $\ge c_j\overline d_j$ when $c_j<0$; summing gives the first
term of \eqref{eq:lbx}. (b) $\partial h/\partial x_d=\sum_jc_j\sig''(z_j)W_{jd}$ and $|\sig''(z_j(x))|\le\kappa_j$ on
$Q$ by the unimodality of $|\sig''|$ on each half-line (if the interval contains $\pm t^\ast$ the bound is the global
maximum; otherwise $|\sig''|$ is monotone on each part of the interval on either side of $0$, where it vanishes, and the
maximum is at an endpoint). The mean value theorem gives $h(x)\ge h(m)-\sum_d|x_d-m_d|\sup_Q|\partial_dh|\ge
h(m)-\sum_dr_d\sum_j|c_j|\kappa_j|W_{jd}|$. (c) $\varphi$ is a sum of functions of one coordinate each, so its minimum
over a product of intervals is the sum of the one-dimensional minima, which are given by the unimodality of $\sig'$ as
in (a).

\paragraph{\Cref{thm:sound}}
Every point of $\B$ belongs to a box of the initial partition, and every box is either split into two boxes that cover it
or discarded. A box is discarded only if $L(Q)\ge0$, which by the bounds implies $h\ge0$ on $Q$; or, in pre-activation
space, if $\min_R\varphi\ge0$, which implies $h(x)=\varphi(Wx+b)\ge0$ for every $x$ with $Wx+b\in R$; or if $R$ does not
meet $Z$, in which case no $x\in\B$ maps into $R$. If the search ends with an empty list, every point of $\B$ lies in a
discarded region where $h\ge0$, so $\min_\B h_i\ge0$ and, by \cref{lem:gain}, $s_ig_i\ge0$ on $\B$. A counterexample is
returned only after evaluating $h(\hat x)<0$ at a point $\hat x\in\B$ (the centre of a sub-box of $\B$, or the solution of
a linear program with bounds $0\le x\le1$).

\paragraph{\Cref{prop:term}}
Inputs with $\sum_j|c_j||W_{jd}|=0$ do not affect $h$ and are never split; along every other coordinate the splitting
rule halves the largest weighted width, so the radius of every box at depth $k$ tends to zero with $k$. Let
$G=\sum_d\sum_j|c_j|(\sqrt3/18)|W_{jd}|$. If $\min_\B h=\delta>0$, the mean-value term of \eqref{eq:lbx} is at least
$\delta-G\max_dr_d$, positive once $\max_dr_d<\delta/G$, so every branch is discarded at finite depth. If
$\min_\B h<0$, continuity gives a box $Q^\ast\subseteq\B$ on which $h<0$; boxes containing a point of $Q^\ast$ are never
discarded (their lower bound is negative), and at finite depth one of them lies inside $Q^\ast$, whose centre is then a
counterexample. In both cases the tree is finite.

\section{Gradients of the gains}\label{app:gradients}
Write $a_j=\sig(z_j)$, $\sig_j'=a_j(1-a_j)$, $\sig_j''=\sig_j'(1-2a_j)$, $o=\sig(z_o)$, $\sig_o'=o(1-o)$,
$\sig_o''=\sig_o'(1-2o)$ and $u_i=\sum_jv_j\sig_j'W_{ji}$, so that $g_i=\sig_o'u_i$. Then, for $k=1,\dots,n$ and
$l=1,\dots,p$,
\begin{align*}
\frac{\partial g_i}{\partial W_{kl}}&=\sig_o''\,v_k\sig_k'x_l\,u_i+\sig_o'\bigl(v_k\sig_k''x_lW_{ki}+v_k\sig_k'\,[i=l]\bigr),\\
\frac{\partial g_i}{\partial b_k}&=\sig_o''\,v_k\sig_k'\,u_i+\sig_o'\,v_k\sig_k''W_{ki},\\
\frac{\partial g_i}{\partial v_k}&=\sig_o''\,a_k\,u_i+\sig_o'\,\sig_k'W_{ki},\\
\frac{\partial g_i}{\partial c}&=\sig_o''\,u_i ,
\end{align*}
where $[i=l]$ is 1 if $i=l$ and 0 otherwise.
These expressions are checked against finite differences in the test suite of the code.

\section{Signal vectors}\label{app:signals}
\begin{table}[H]\centering\small
\caption{Expected gain signs. $+$: $P(Y=1\mid x)$ non-decreasing in the input; $-$: non-increasing; $0$: unconstrained.
The signs of the five original datasets were derived from domain knowledge (supplementary material); those of the
benchmarks are the ones used in the literature \citep{Liu2020,Sivaraman2020,Runje2023}.}\label{tab:signals}
\begin{tabular}{lp{0.8\textwidth}}\toprule
ID & Inputs (sign) \\\midrule
ALG & temperature ($+$), relative humidity ($-$), wind speed (0), rain ($-$), FFMC, DMC, DC, ISI, BUI, FWI ($+$) \\
BCO & clump thickness, uniformity of cell size, uniformity of cell shape, marginal adhesion, single epithelial cell size,
bare nuclei, bland chromatin, normal nucleoli, mitoses (all $+$) \\
BCD & mean, standard error and worst value of radius, texture, perimeter, area, smoothness, compactness, concavity,
concave points and symmetry ($+$); fractal dimension (0) \\
HF & age ($+$), anaemia ($+$), creatinine phosphokinase (0), diabetes (0), ejection fraction ($-$), high blood pressure
($+$), platelets (0), serum creatinine ($+$), serum sodium ($-$), sex (0), smoking (0), follow-up time ($-$) \\
PIM & pregnancies ($+$), glucose ($+$), blood pressure (0), skin thickness ($+$), insulin (0), BMI ($+$), diabetes
pedigree function ($+$), age ($+$) \\
HD & resting blood pressure ($+$), serum cholesterol ($+$); the other eleven inputs (0) \\
COM & prior adult convictions, juvenile felonies, juvenile misdemeanours, other juvenile counts ($+$); age, race and sex
indicators (0) \\
LOAN & credit score, length of employment, annual income ($-$); public-record bankruptcies, debt-to-income ratio ($+$);
23 other inputs (0) \todo{confirm the order of the first five features in the Zenodo release} \\
\bottomrule\end{tabular}
\end{table}

\section{Hyper-parameter grids}\label{app:hyper}
\begin{table}[H]\centering\small
\caption{Hyper-parameter grids of the baselines.}\label{tab:hyper}
\resizebox{\textwidth}{!}{\begin{tabular}{ll}\toprule
Method & Grid (selected on the validation split) \\\midrule
LR, LR-C & $L_2$ penalty $\in\{10^{-4},10^{-3},10^{-2},10^{-1}\}$; LR-C: coefficient signs bounded by $s$ \\
XGB(-C) & depth $\in\{2,3,4\}$, trees $\in\{100,300\}$, learning rate 0.05, subsample 0.9 \\
LGBM(-C) & leaves $\in\{4,8,16\}$, trees $\in\{100,300\}$, learning rate 0.05, subsample 0.9 \\
MLP & one hidden layer of $\{1,2,4,\dots,64\}$ logistic units \\
Min-Max & $K$ groups of $K$ hyperplanes, $K\in\{2,4,8\}$ \\
CMNN & width $\in\{8,16,32\}$, two monotone layers, ELU-based activations \\
LMN & width $\in\{16,32\}$, Lipschitz constant $\lambda\in\{2,8,32\}$, GroupSort \\
Min-Max, CMNN, LMN & Adam, learning rate $5\times10^{-3}$, batch 256, early stopping (patience 50, at most 1000 epochs) \\
\bottomrule\end{tabular}}
\end{table}

\bibliographystyle{elsarticle-harv}
\bibliography{refs}

\end{document}
